\lampiransection{INDEKS KINERJA NASI GERILYA DAN KOMPETITOR PADA GRABFOOD}
\label{app:prasurvey3}

Lampiran ini memuat hasil pra-survei tahap 3 penelitian \textit{Analisis Strategi Pemasaran Rumah Makan Nasi Gerilya pada Platform GrabFood Menggunakan Matriks SWOT}. Dokumen sumber berasal dari pihak Nasi Gerilya dalam bentuk dokumentasi internal performa toko dan pembanding \textit{merchant} pada platform GrabFood. Karena nominal \textit{revenue} eksak terikat perjanjian kerahasiaan (\textit{non-disclosure agreement}/NDA), penyajian pada lampiran ini dibatasi dalam bentuk indeks turunan. Indeks omzet disusun dengan basis Desember 2024 = 100, indeks AOV dengan basis Juli 2025 = 100, indeks kunjungan harian dengan basis September 2025 = 100, dan indeks kompetitor dengan basis Juli 2025 = 100 untuk masing-masing \textit{merchant}.

\begin{table}[h!]
    \centering
    \caption{Indeks Omzet Bulanan Nasi Gerilya}
    \label{tab:prasurvey3_indeks_omzet}
    \begin{tabular}{lcc}
        \toprule
        \textbf{Periode} & \textbf{Indeks} & \textbf{Perubahan} \\
        \midrule
        Des-24           & 100,0           & --                 \\
        Jan-25           & 103,0           & +3,0\%             \\
        Feb-25           & 112,6           & +9,3\%             \\
        Mar-25           & 99,1            & -12,0\%            \\
        Apr-25           & 116,0           & +17,0\%            \\
        Mei-25           & 175,3           & +51,1\%            \\
        Jun-25           & 383,8           & +118,9\%           \\
        Jul-25           & 633,3           & +65,0\%            \\
        Ags-25           & 104,2           & -83,6\%            \\
        Sep-25           & 103,2           & -1,0\%             \\
        Okt-25           & 93,8            & -9,1\%             \\
        Nov-25           & 86,5            & -7,8\%             \\
        \bottomrule
    \end{tabular}
\end{table}

\begin{table}[h!]
    \centering
    \caption{Indeks \textit{Average Order Value} (AOV)}
    \label{tab:prasurvey3_indeks_aov}
    \begin{tabular}{lcc}
        \toprule
        \textbf{Periode} & \textbf{Indeks} & \textbf{Perubahan} \\
        \midrule
        Jul-25           & 100,0           & --                 \\
        Ags-25           & 102,0           & +2,0\%             \\
        Sep-25           & 107,9           & +5,8\%             \\
        Okt-25           & 110,0           & +1,9\%             \\
        Nov-25           & 114,5           & +4,1\%             \\
        \bottomrule
    \end{tabular}
\end{table}

\begin{table}[h!]
    \centering
    \caption{Indeks Kunjungan Harian}
    \label{tab:prasurvey3_indeks_kunjungan}
    \begin{tabular}{lcc}
        \toprule
        \textbf{Periode} & \textbf{Indeks} & \textbf{Perubahan} \\
        \midrule
        Sep-25           & 100,0           & --                 \\
        Okt-25           & 116,5           & +16,5\%            \\
        Nov-25           & 130,7           & +12,2\%            \\
        \bottomrule
    \end{tabular}
\end{table}

\begin{table}[h!]
    \centering
    \small
    \caption{Indeks Kompetitor Nasi Padang Juli--Oktober 2025}
    \label{tab:prasurvey3_indeks_kompetitor}
    \begin{tabularx}{\textwidth}{p{4cm}ccccc}
        \toprule
        \textbf{\textit{Merchant}} & \textbf{Jul} & \textbf{Ags} & \textbf{Sep} & \textbf{Okt} & \textbf{$\Delta$ Jul--Okt} \\
        \midrule
        Koto Gadang                & 100,0        & 105,1        & 86,2         & 94,2         & -5,8\%                     \\
        Pondok Krakatau            & 100,0        & 110,1        & 89,0         & 107,5        & +7,5\%                     \\
        Uda Sayang                 & 100,0        & 94,1         & 89,0         & 88,9         & -11,1\%                    \\
        Dendeng Batokok            & 100,0        & 104,1        & 85,2         & 95,1         & -4,9\%                     \\
        Istana Krakatau            & 100,0        & 119,2        & 108,4        & 119,3        & +19,3\%                    \\
        Sederhana                  & 100,0        & 97,5         & 94,9         & 98,6         & -1,4\%                     \\
        Gumarang Jaya              & 100,0        & 116,1        & 102,6        & 114,1        & +14,1\%                    \\
        Nasi Gerilya               & 100,0        & 92,6         & 87,8         & 70,4         & -29,6\%                    \\
        Garuda                     & 100,0        & 95,5         & 90,1         & 93,9         & -6,1\%                     \\
        Pondok Gurih               & 100,0        & 102,9        & 97,0         & 101,5        & +1,5\%                     \\
        \bottomrule
    \end{tabularx}
\end{table}

Data pada Tabel~\ref{tab:prasurvey3_indeks_omzet}, \ref{tab:prasurvey3_indeks_aov}, \ref{tab:prasurvey3_indeks_kunjungan}, dan \ref{tab:prasurvey3_indeks_kompetitor} digunakan sebagai dasar pembacaan arah performa usaha tanpa membuka nominal \textit{revenue} absolut yang bersifat rahasia.
